\documentclass[10pt,a4paper]{amsart}
\usepackage{amsmath,amssymb,amsthm,graphicx}
\usepackage{verbatim}
\newenvironment{reference}{\comment}{\endcomment}
%\newenvironment{reference}{}{}
\newenvironment{slogan}{\comment}{\endcomment}
\newenvironment{history}{\comment}{\endcomment}

% For commutative diagrams we use Xy-pic
\usepackage[all]{xy}

% We use 2cell for 2-commutative diagrams.
\xyoption{2cell}
% Theorem environments.
%
\theoremstyle{plain}
\newtheorem{theorem}[subsection]{Theorem}
\newtheorem{proposition}[subsection]{Proposition}
\newtheorem{lemma}[subsection]{Lemma}

\theoremstyle{definition}
\newtheorem{definition}[subsection]{Definition}
\newtheorem{example}[subsection]{Example}
\newtheorem{exercise}[subsection]{Exercise}
\newtheorem{situation}[subsection]{Situation}

\theoremstyle{remark}
\newtheorem{remark}[subsection]{Remark}
\newtheorem{remarks}[subsection]{Remarks}

\numberwithin{equation}{subsection}

% Macros
%
\def\lim{\mathop{\mathrm{lim}}\nolimits}
\def\colim{\mathop{\mathrm{colim}}\nolimits}
\def\Spec{\mathop{\mathrm{Spec}}}
\def\Hom{\mathop{\mathrm{Hom}}\nolimits}
\def\Ext{\mathop{\mathrm{Ext}}\nolimits}
\def\SheafHom{\mathop{\mathcal{H}\!\mathit{om}}\nolimits}
\def\SheafExt{\mathop{\mathcal{E}\!\mathit{xt}}\nolimits}
\def\Sch{\mathit{Sch}}
\def\Mor{\mathop{\mathrm{Mor}}\nolimits}
\def\Ob{\mathop{\mathrm{Ob}}\nolimits}
\def\Sh{\mathop{\mathit{Sh}}\nolimits}
\def\NL{\mathop{N\!L}\nolimits}
\def\CH{\mathop{\mathrm{CH}}\nolimits}
\def\proetale{{pro\text{-}\acute{e}tale}}
\def\etale{{\acute{e}tale}}
\def\QCoh{\mathit{QCoh}}
\def\Ker{\mathop{\mathrm{Ker}}}
\def\Im{\mathop{\mathrm{Im}}}
\def\Coker{\mathop{\mathrm{Coker}}}
\def\Coim{\mathop{\mathrm{Coim}}}

% Boxtimes
%
\DeclareMathSymbol{\boxtimes}{\mathbin}{AMSa}{"02}

%
% Macros for moduli stacks/spaces
%
\def\QCohstack{\mathcal{QC}\!\mathit{oh}}
\def\Cohstack{\mathcal{C}\!\mathit{oh}}
\def\Spacesstack{\mathcal{S}\!\mathit{paces}}
\def\Quotfunctor{\mathrm{Quot}}
\def\Hilbfunctor{\mathrm{Hilb}}
\def\Curvesstack{\mathcal{C}\!\mathit{urves}}
\def\Polarizedstack{\mathcal{P}\!\mathit{olarized}}
\def\Complexesstack{\mathcal{C}\!\mathit{omplexes}}
% \Pic is the operator that assigns to X its picard group, usage \Pic(X)
% \Picardstack_{X/B} denotes the Picard stack of X over B
% \Picardfunctor_{X/B} denotes the Picard functor of X over B
\def\Pic{\mathop{\mathrm{Pic}}\nolimits}
\def\Picardstack{\mathcal{P}\!\mathit{ic}}
\def\Picardfunctor{\mathrm{Pic}}
\def\Deformationcategory{\mathcal{D}\!\mathit{ef}}
\title{An affine scheme and an affine open whose complement admits no quasi-compact dense open}
\author{The Stacks Project Authors}
\date{stacks-project:examples.tex:local-snapshot-20261012; GNU FDL 1.2 or later}
\begin{document}
\maketitle
\noindent\textit{Editorial note.} All delivered mathematical text is the author's own, copied from the pinned Stacks Project chapter file examples.tex; no statement or proof step was rewritten, reordered or removed.\par\medskip
\noindent\textit{Editorial note.} The counted proof is the author's complete argument: The statement's item 1 is witnessed by (T, U) and item 2 by (S, s); the (T, t) = (S, s) isomorphism used later in the chapter is a separate remark, not needed here.\par\medskip
\noindent\textit{Editorial note.} The label lemma-complement-of-affine-does-not-contain-qc-dense-open names the counted statement, and the proof environment immediately following it in the source is delivered verbatim as its proof.\par\medskip
\noindent\textit{Editorial note.} The counted statement is carried by the author's own lemma environment in the chapter source.\par\medskip
\noindent\textit{Allowed standard prerequisites (editorial).}\begin{itemize}
\item Context span 2249-2261: A = k[z\_1, z\_2, ...]/(all pairwise products z\_i z\_j, i != j), S = Spec(A), s the closed point of the maximal ideal (z\_i); S - \{s\} = union of the D(z\_i), each D(z\_i) open irreducible with generic point eta\_i; any quasi-compact open V of S - \{s\} dense in S - \{s\} would contain all eta\_i, but a principal affine open D(f) with f in (z\_1, z\_2, ...) lies in (z\_1, ..., z\_n) for some n and contains only finitely many eta\_i, so V cannot be quasi-compact
\item Context span 2263-2274: B = k[x, z\_1, z\_2, ...]/(x z\_i for all i and all pairwise products z\_i z\_j), T = Spec(B), U = D(x); for t the closed point of (x, z\_i) the closure of U is U cup \{t\}; a quasi-compact open V with V cap U = empty and U cup V dense in T would be a quasi-compact open of union\_i D(z\_i) dense in union\_i D(z\_i), which the first paragraph excludes
\end{itemize}\medskip
Let $k$ be a field and let $A = k[z_1, z_2, z_3, \ldots]/I$ where $I$
is the ideal generated by all pairwise products
$z_iz_j$, $i \not = j$, $i, j \in \mathbf{N}$. Set $S = \Spec(A)$.
Let $s \in S$ be the closed point corresponding to the maximal
ideal $(z_i)$. We claim there is no
quasi-compact open $V \subset S \setminus \{s\}$ which is
dense in $S \setminus \{s\}$. Note that $S \setminus \{s\} = \bigcup D(z_i)$.
Each $D(z_i)$ is open and irreducible with generic point
$\eta_i$. We conclude that $\eta_i \in V$ for all $i$.
However, a principal affine open of $S \setminus \{s\}$ is
of the form $D(f)$ where $f \in (z_1, z_2, \ldots)$. Then
$f \in (z_1, \ldots, z_n)$ for some $n$ and we see that $D(f)$ contains
only finitely many of the points $\eta_i$. Thus $V$ cannot be quasi-compact.

\medskip\noindent
Let $k$ be a field and let $B = k[x, z_1, z_2, z_3, \ldots]/J$ where $J$
is the ideal generated by the products $xz_i$, $i \in \mathbf{N}$ and by
all pairwise products $z_iz_j$, $i \not = j$, $i, j \in \mathbf{N}$.
Set $T = \Spec(B)$. Consider the principal open $U = D(x)$.
We claim there is no quasi-compact open $V \subset T$ such that
$V \cap U = \emptyset$ and $V \cup U$ is dense in $T$.
Let $t \in T$ be the closed point corresponding to the maximal
ideal $(x, z_i)$. The closure of $U$ in $T$ is
$\overline{U} = U \cup \{t\}$. Hence $V \subset \bigcup_i D(z_i)$
is a quasi-compact open. By the arguments of the previous paragraph
we see that $V$ cannot be dense in $\bigcup D(z_i)$.

\begin{lemma}
\label{lemma-complement-of-affine-does-not-contain-qc-dense-open}
Nonexistence quasi-compact opens of affines:
\begin{enumerate}
\item There exist an affine scheme $S$ and affine open $U \subset S$
such that there is no quasi-compact open $V \subset S$ with
$U \cap V = \emptyset$ and $U \cup V$ dense in $S$.
\item There exists an affine scheme $S$ and a closed point $s \in S$ such that
$S \setminus \{s\}$ does not contain a quasi-compact dense open.
\end{enumerate}
\end{lemma}

\begin{proof}
See discussion above.
\end{proof}

\end{document}
